% Insert as A.60.  This supplies all old-place corrections for weight four.
\begin{statement}{A.60}{Proved}{Exact weight-three matching at the places needed for the fourth layer}
The restriction map
\[
 \begin{split}
 H^1_{\rm cts}(G_\Q,\widehat\Z(3))\longrightarrow{}&
 H^1(\Q_2,\Z_7(3))\oplus H^1(\Q_3,\Z_{13}(3))\\
 &\oplus H^1(\Q_3,\Z_2(3))\oplus H^1(\R,\Z_2(3))
 \end{split}
 \tag{A.60.1}
\]
is surjective.  The four target groups have orders
$7,13,2,2$, respectively.

In particular, fix an ordinary pair with valuation patterns
$(0,1)$ at two and $(1,0)$ at three, and an actual
class-three lift of its distinguished class-two section $s_*$.
There is an actual class-three lift above the same $s_*$
which agrees, up to geometric conjugacy, with the ordinary
point sections at the following local quotient maps:
\[
 (q,\ell)=(2,7),\ (3,13),\ (3,2),\ (\infty,2).
 \tag{A.60.2}
\]
The quotient map here is the projection of the geometric
nilpotent group from its $\widehat\Z$ Hall coordinates to
its $\Z_\ell$ Hall coordinates; it preserves the arithmetic
Galois group.  The two weight-three central coordinates
are corrected independently by (A.60.1).
\end{statement}
\begin{proof}
For $q\ne\ell$, a continuous integral weight-three
one-cocycle is unramified: its tame inertia value $t$
satisfies $q^3t=qt$, and $\Z_\ell$ is torsion-free.
Wild inertia has no nontrivial homomorphism into this
pro-$\ell$ additive module.  Consequently
\[
 H^1(\Q_q,\Z_\ell(3))
       =H^1_{\rm ur}(\Q_q,\Z_\ell(3))
       =\Z_\ell/(q^3-1)\Z_\ell.
 \tag{A.60.3}
\]
At the real place the odd twist gives
$H^1(\R,\Z_2(3))=\Z_2/2\Z_2$.
These formulas give the four orders in the statement.

\emph{The two odd coefficient primes.}
For $\ell=7$ or $13$, let $F=\Q(\zeta_\ell)$ and set
\[
              u_\ell=\prod_{a=1}^{\ell-1}
                            (1-\zeta_\ell^a)^{a^2}.
 \tag{A.60.4}
\]
This is an $S$-unit of $F$, with $S$ the places above
$\ell$ and the infinite places.  Under
$\sigma_b(\zeta_\ell)=\zeta_\ell^b$, reindexing the
product shows
\[
       \sigma_b(u_\ell)=u_\ell^{b^{-2}}
                       \quad\hbox{in }F^\times/(F^\times)^\ell.
 \tag{A.60.5}
\]
Indeed the reindexed exponents differ from
$b^{-2}a^2$ by multiples of $\ell$.
Twisting the Kummer class in $H^1(F,\F_\ell(1))$
by $\F_\ell(2)$ therefore makes it invariant under
$\Delta=\operatorname{Gal}(F/\Q)$.
Because $|\Delta|=\ell-1$ is invertible on these
coefficients, inflation--restriction identifies this
invariant class with a class
\[
          \bar c_\ell\in
          H^1(G_{\Q,\{\ell,\infty\}},\F_\ell(3)).
 \tag{A.60.6}
\]
The Kummer extension is unramified outside $\ell$:
away from $\ell$, its defining unit and its derivative
$\ell T^{\ell-1}$ are both invertible.
By the vanishing
$H^2_{\rm cts}(G_{\Q,\{\ell,\infty\}},\Z_\ell(3))=0$
proved in A.46, the coefficient connecting map is zero,
so $\bar c_\ell$ lifts to an integral class $c_\ell$.

The needed local values are exact finite-field computations.
Use
\[
 \mathbf F_8=\mathbf F_2[z]/(z^3-z-1),\qquad
 \mathbf F_{27}=\mathbf F_3[z]/(z^3-z-1).
 \tag{A.60.7}
\]
In the first field $z$ has order seven, and in the second
it has order thirteen.  Direct multiplication of
(A.60.4), with $\zeta_\ell$ reduced to $z$, gives
\[
 \begin{array}{c|c|c}
 (\ell,q)&u_\ell\bmod q&
          u_\ell^{(q^3-1)/\ell}\bmod q\\\hline
 (7,2)&z&z,\\
 (13,3)&1+z&z^6.
 \end{array}
 \tag{A.60.8}
\]
The irreducible polynomials, element orders, products,
and displayed identities are checked by exact finite-field
arithmetic in \texttt{weight3\_local\_generators.py}.
Both residue degrees in $F/\Q$ are three.
The weight-three action of rational Frobenius is one
modulo $\ell$ at these primes, so restriction to its
third power multiplies its cocycle value by three.
The descended Frobenius values are consequently
$3^{-1}=5\pmod7$ and $6\cdot3^{-1}=2\pmod{13}$.
They are nonzero.  Reduction on the two integral local
groups in (A.60.3) is injective, and those groups have
prime order.  Thus $c_7$ and $c_{13}$ generate the first
two target factors.

\emph{The real and dyadic factors.}
A.46 also gives
$H^2_{\rm cts}(G_{\Q,\{2,\infty\}},\Z_2(3))=0$.
Hence both finite classes $\kappa_2(-1)$ and
$\kappa_2(2)$ lift to integral weight-three classes
unramified outside two.  In the order
$(\R,\Q_3)$ their reductions have restriction vectors
\[
                         (1,1),\qquad(0,1).
 \tag{A.60.9}
\]
At the real place these are their signs.  At three
both classes are unramified, and both $-1$ and $2$
are nonsquares in $\mathbf F_3^\times$.
Reduction identifies the two integral groups in question
with these order-two finite images.  The matrix of
(A.60.9) has determinant one over $\F_2$.
Thus the last two factors can be prescribed independently.
Taking the three coefficient-prime components in the
product of continuous cochain complexes proves (A.60.1).

Finally, at $(q,\ell)=(2,7)$ and $(3,13)$, the
ordinary abelian point coordinates already determine the
class-two section: its central comparison group is
$\Z_\ell/(q^2-1)\Z_\ell=0$.
At $(3,2)$ this agreement is the defining normalization
of $s_*$, and at the real place
$H^1(\R,\Z_2(2))=0$ gives the same agreement.
Choose geometric conjugations aligning these local
class-two representatives, and lift those conjugators
to class three.  The resulting differences of class-three
sections are central one-cocycle classes with coefficients
$\Z_\ell(3)^2$.  This is the actual fibre action of the
extension with kernel $\Gamma_3/\Gamma_4$ from A.43.
Apply (A.60.1) to each coordinate and subtract the
resulting global central cocycle from the chosen
class-three lift.  Its image in $E_2$ is unchanged,
and its four local class-three classes are the desired
point classes.  At the real place any chosen alignment
suffices; uniqueness of that alignment is not asserted.
\end{proof}
